\section{Protobuf Compatibility Relations}~\label{sec:comp-rel}

Compatibility between protobuf entities is defined with a series of
relations. The first and most self-contained relation is the value relation,
which relates two values at the specification level if serializing just a value
(such as a variable width integer) and then parsing it to a different protobuf
type would lead to a different value. After that is the protobuf type relation,
which relates two protobuf types if any value of the first type is related to
some value of the second type in the value relation. Since protobuf field and
message definitions will not have concrete values associated with them, the type
relation can be used in place of the value relation to make safe type
changes. Finally, the message relation connects compatible messages such that
all instances of the original message can be parsed into instances of the updated
message.

Below is a list of all meta-variables used across all relations, and even across
encoding formats. When multiple of these variables are needed, you will see
variants with subscript numbers or apostrophes.

\begin{align*}
  m &: \text{message} & v &: \text{value} \\
  id &: \text{field number} & \tau &: \text{type} \\
  r &: \text{Reserved field set} & \mathtt{n},\ \mathtt{m} &: \text{bit width} \\
  f &: \text{Field map} & e &: \text{enum} \\
  s &: \text{Field name} & d &: \text{decorator}
\end{align*}

Below is a list of the relations about to be defined, along with an English
explanation of how to read them.

\begin{align*}
  \prec &: \text{Value relation} & v_1 : \tau_1 \prec v_2 : \tau_2 &\rightarrow ``v_1 : \tau_1 \text{ becomes } v_2 :
                                           \tau_2 " \\
  \propto &: \text{Type relation} & \tau_1 \propto \tau_2 &\rightarrow ``\tau_1\text{ is type-compatible with }
                                          \tau_2" \\
  \preceq &: \text{Message relation} & m_1 \preceq m_2 &\rightarrow ``m_1 \text{ is compatible with } m_2"
\end{align*}


\subsection{Value Relation}~\label{sec:val-rel}

The value relation is the base level of the compatibility relations, relating
just two individual values. Protobuf values are modeled here as a specification
level type in Lean associated with a protobuf type as listed in
Table~\ref{tab:val-spec}. Finally, while protobuf decorators are technically
part of the field level specification, they have been incorporated into the type
information to allow the specification type into include \texttt{option} or
\texttt{list} Lean types.

\begin{table}[H]
	\centering
	\begin{tabular}{cl}
      \toprule
      Specification Type & Protobuf Type \\
      \midrule
      $\mathbb{Z}$ & \ints, \uints, \sints,
            \intl, \uintl, \sintl{} \\
      $\mathbb{B}$ & \bool{}    \\
      \str{} & \str{}    \\
      \texttt{list w8}  & \byt{}    \\
      Bounded Integer & \texttt{float}, \texttt{double} \\
      \bottomrule
	\end{tabular}

	\vspace{4mm}
	\caption[]{Description of the Lean type corresponding to each protobuf
      type. Floating point numbers are included in the model, but don't have a
      meaningful compatibility story, so they don't appear in the relations
      defined below.}\label{tab:val-spec}
\end{table}

While it might be tempting to have a roundtrip theorem where the output of the
serializer is passed directly to the parser, this would require a canonical
encoding, which Protobuf doesn't have. Instead we make use of a relational
encoding specification, where $E_{\tau}(v, bs)$ means that $v$ can be encoded
as bytes $bs$ under type $\tau$. \\

\begin{definition}[Value Relation]
  A protobuf value at the specification level $v_1$ with protobuf type $\tau_1$ is
  related to another value $v_2$ at type $\tau_2$, denoted $v_1 : \tau_1 \prec v_2 : \tau_2$
  if $\forall\ bs, E_{\tau}(v_1, bs) \rightarrow parse_{\tau_2}(bs) = v_2$.
\end{definition}

\subsubsection{Basic Rules}

The value relation is reflexive and transitive, however a direct transitivity
rules isn't included in the relation. To understand why, first consult the
integer rules in Section~\ref{sec:proto-compat-ints} and consider what happens
when an integer is cast to a smaller byte width. Specifically, consider $2^{40}$
as a \texttt{uint64}. Casting this value to a \texttt{uint32} can be done with
the \textsc{Uint-Chg-W} rule, resulting in $2^{40} : \mathtt{uint64} \prec 0 :
\mathtt{uint32}$, which describes changing the field perfectly. However, we can
then promote this field back to \texttt{uint64} with another application of
\textsc{Uint-Chg-W}, getting $0 : \mathtt{uint32} \prec 0 :
\mathtt{uint64}$. Having a transitive rule would result in $2^{40} :
\mathtt{uint64} \prec 0 : \mathtt{uint64}$ but clearly parsing a serialization
of $2^{40}$ as a \texttt{uint64} when it was serialized as a \texttt{uint64}
should recover the original value!

\begin{mathpar}
  \infer[Refl]{ }{v:\tau \prec v:\tau}

  % \infer[Trans]{v_1:\tau_1 \prec v_2:\tau_2 \\ v_2:\tau_2 \prec v_3:\tau_3}{v_1:\tau_1 \prec v_3:\tau_3}
\end{mathpar}

\subsubsection{String \& Byte Rules}

It is possible to convert been the \texttt{string} and \texttt{bytes} types,
although the host program may have to deal with control character bytes being
present in the resulting string. Technically the Protobuf specification requires
that all strings are encoded in UTF-8, so we can optionally add that requirement
to converting from \texttt{bytes} to \texttt{string}.

\begin{mathpar}
  \infer[Str-Byte]{ }{v: \str{} \prec{} v: \byt{}}

  \infer[Byte-Str]{ {\color{red} v \text{ is valid UTF-8}} }{v: \byt{} \prec{} v: \str{}}
\end{mathpar}

\subsubsection{Varint Scalar Rules}\label{sec:proto-compat-ints}

Every integer type and several non-integer protobuf types are encoded using the
variable-width integer encoding described in Section~\ref{sec:proto-vint}. To
compactly represent 30 rules, we'll observe that all the varint-encoded values
are transmitted over the wire in the range $\payload$. We can factor each
rule as using the writer's type to convert it to the wire value, then the
reader's type to convert it to the new parsed value.

As a shorthand, statements like \uintn{n} represent a variable length integer
encoding into $n$ bits. In order to result in valid protobuf types,
$n \in \{32, 64\}$. 

It is also worth noting that the expression $v : \tau$ actually refers to a value
of type $\denote{\tau}$ with a label of the type. All of the integer
types have $\denote{\cdot} = \mathbb{Z}$, as shown in Table~\ref{tab:val-spec},
which allows for the free conversion between integer types. Finally, $\%$ refers
to modulus using floored division, so $-1 \mod 2 = 1$ and the rules are
written with truncated integer division.

Let $\vints = \left\{ \uintn{n}, \intn{n}, \sintn{n}, \bool \right\}$ represent
the protobuf types encoded on the wire as a variable-width integer. Some rules
might show enums, which are encoded this way, but until their formalization is
planned, these are all draft rules. For each $\tau \in \vints$, let
\begin{align*}
  \encop{\tau} &: \denote{\tau} \rightarrow \payload \\
  \decop{\tau} &: \payload \rightarrow \denote{\tau}
\end{align*}
be functions that convert from a $\tau$-typed writer to the payload number and
the value recovered by a $\tau$-typed reader respectively.
\[
  \infer[Varint]{ \tau_1, \tau_2 \in \vints \\ v_2 =
    \dec{\tau_2}{\enc{\tau_1}{v_1}}}{v_1 : \tau_1 \prec v_2 : \tau_2}
\]

The table below gives the decoding and encoding functions for each type, and
makes use of the following definitions.

\begin{table}[H]
  \centering
  \begin{tabular}{l c c}
    \toprule
    $\tau$ & $\enc{\tau}{v}$ & $\dec{\tau}{x}$ \\
    \midrule
    \texttt{uint[n]} & $v$ & $x \mod 2^n$ \\
    \texttt{int[n]} & $v \mod 2^{64}$ & $((x + 2^{n-1}) \mod 2^n) - 2^{n-1}$ \\
    \texttt{sint[n]} & $\zz{v}$ & $\unzz{x \mod 2^n}$ \\
    \texttt{bool} & $\ind{v}$ & $x \ne 0$ \\
    {\color{red} \texttt{enum} $e$} & $v \mod 2^{64}$ & $((x + 2^{31}) \mod 2^{31}) - 2^{31}$ \\
    \bottomrule
  \end{tabular}
  \caption[Encoding and decoding rules for Varint conversions.]{Payload written
    and value recovered for each varint-encoded type. Here $\%$ is always the
    non-negative floored modulus; no truncated division occurs. The
    $\encop{\intn{n}}$ row is the ten-byte sign extension: protobuf writes every
    negative $\intn{n}$ as its two's complement at 64 bits, which is exactly
    $v \mod 2^{64}$.
  }
  \label{tab:varint-conversion}
\end{table}

Which uses these definitions for converting to and from the zig-zag encoding
(See Section~\ref{sec:proto}).
\begin{align*}
  \zz{v} &= \left\{
    \begin{array}{l@{\quad:\quad}l}
      2v & v \ge 0 \\
      -2v - 1 & v < 0
    \end{array}\right. \\
  \unzz{x} &= \left\{
    \begin{array}{l@{\quad:\quad}l}
      x / 2 & x \text{ even} \\
      -(x + 1) / 2 & x \text{ odd.}
    \end{array}\right.
\end{align*}

Using these functions, we can still derive and simplify every rule for
converting between varint scalars.

\begin{table}[H]
  \centering
  \begin{tabular}{lcccc}
    \toprule
    \multirow{2}{*}{writer $v_1 : \tau_1$} & \multicolumn{4}{c}{reader $\tau_2$} \\
    & \uintn{m} & \intn{m} & \sintn{m} & \bool{} \\
    \midrule
    \uintn{n} & $v_1 \mod 2^m$      & $\sigma_m(v_1)$      & $\unzz{v_1 \mod 2^m}$      & $v_1 \ne 0$ \\
    \intn{n}  & $v_1 \mod 2^m$      & $\sigma_m(v_1)$      & $\unzz{v_1 \mod 2^m}$      & $v_1 \ne 0$ \\
    \sintn{n} & $\zz{v_1} \mod 2^m$ & $\sigma_m(\zz{v_1})$ & $\unzz{\zz{v_1} \mod 2^m}$ & $v_1 \ne 0$ \\
    \bool{}   & $\ind{v_1}$   & $\ind{v_1}$    & $-\ind{v_1}$         & $v_1$ \\
    \bottomrule
  \end{tabular}
  \caption{The transform induced by \textsc{Varint} on every pair of
    varint-encoded types, writing
    $\sigma_m(x) = \left((x + 2^{m-1}) \mod 2^m\right) - 2^{m-1}$ for the
    reader's signed truncation.  Widths $n$ and $m$ are independent, so this
    covers all thirty non-identity conversions.}\label{tab:varint-derived}
\end{table}

\subsubsection{Message \& Enum Rules}

Definitions for the structure of a message and enum are given in
Section~\ref{sec:msg-rel}. The values are represented as a mapping from the id
number of each field to the value and type of that field. 

\begin{mathpar}
  \infer[Msg]{\msg{r_1}{f_1} \preceq \msg{r_2}{f_2} \\ \forall\;i \in \dom{v_2}.\;\left( f_1(i) =
      (s, \tau_1) \wedge v_1(i) \prec v_2(i) \right) \vee v_2(i) = default\left(\mathtt{snd}\ f_2(i)\right)}
  {v_1 : \mathtt{MSG}\ \msg{r_1}{f_1} \prec v_2 : \mathtt{MSG}\ \msg{r_2}{f_2} }
  
  \infer[Enum]{v \in e_2}{v : \enum{e_1} \prec v : \enum{e_2}}
\end{mathpar}

The first rule is complex. It stands that a message with some reserved fields
and field definitions can be interpreted as a message with a different set of
reserved fields and field definitions if the two descriptors are compatible and
if for all fields in the resulting value, either (1) it has a compatible value
(2) it has the default value.

The second rule simply states that the value of an enum $e_1$ is also defined in
enum $e_2$.

\subsubsection{Decorator Rules}

Since the specification type can include \texttt{option}'s and
\texttt{repeated}'s, rules for handling these must also be present in the value
relation. These rules are grouped into introduction rules, which take implicit
values and make them optional or repeated, pass through rules which enable the
above rules to operate on these modified types.

Generally speaking, we can introduce, eliminate {\color{red} (It's clear what
  happens when we eliminate a non-packed \texttt{repeated} field; the last
  writer wins. What happens when we eliminate a packed field? We'll end up
  skipping the field as unknown. This difference is actually why we can't have a
  \textsc{Rep-Elim} rule since we can't implement the elimination as a function
  without more knowledge.)} or change the
underlying type of each decorator. Implicit fields require no special
consideration.

{\color{red} There really should be a rule which can introduce a \texttt{None},
  but since implicit fields are modeled as just the spec type, this is
  difficult. I will likely need to model both implicit and optional fields with
  an \texttt{option}, add an annotation to the protobuf type and provide a
  function which introduces a default value for each type. Alternatively, I
  could have a \texttt{MISSING} value and a get function, which if it encounters
  the \texttt{MISSING} just returns the default value. That would align with the
  descriptor setup in my intermediate parse setup.}

\begin{mathpar}
  \infer[Rep-Pass]{v_1:\tau_1 \prec v_1':\tau_2 \\ \cdots \\ v_n:\tau_1 \prec v_n':\tau_2}{[v_1; \dots;
    v_n]: \rept{\tau_1} \prec [v_1'; \dots; v_n']: \rept{\tau_2}}

  \infer[Opt-Some-Pass]{v_1:\tau_1 \prec v_2:\tau_2}{(\some{v_1}): \optt{\tau_1} \prec
    (\some{v_2}): \optt{\tau_2}}

  \infer[Opt-None-Pass]{ }{\none: \optt{\tau_1} \prec \none: \optt{\tau_2}}
  
  \infer[Opt-Intro]{ }{v:\tau \prec (\some{v}): \optt{\tau}}

  \infer[Rep-Intro]{ }{v:\tau \prec [v]: \rept{\tau}}

  \infer[Missing-Imp]{\tau_1 \propto \tau_2}{\none: \tau_1 \prec default(\tau_2):\tau_2}
\end{mathpar}

\subsection{Type Relation}~\label{sec:typ-rel}

This relation relates types which can be converted via the value
relation. Recall that this relation is needed to reason about descriptor-only
compatibility, which doesn't have any concrete values associated with it.

\begin{definition}[Type Relation]
  The type relation relates two protobuf type $\tau_1$ and $\tau_2$, denoted $\tau_1 \propto
  \tau_2$ if for all $v_1:\tau_1$ there exists a $v_2:\tau_2$ such that $v_1:\tau_1 \prec v_2:\tau_2$.
\end{definition}

These rules compress many of the rules in the previous section into one rule,
stating that types are compatible with each other.

\subsubsection{Base Type Rules}

\begin{mathpar}
  \infer[Str-Byt-T]{ }{\str \propto \byt}

  \infer[Byt-Str-T]{ }{\byt \propto \str}

  \infer[Int-Int-T]{\tau_1, \tau_2 \in \left\{{
        \begin{array}{c}
          \ints, \intl, \uints, \uintl \\
          \sints, \sintl, \bool
        \end{array}
      }\right\}}{\tau_1 \propto \tau_2}
\end{mathpar}

It's worth pointing out that converting to and from \sintn{n} is tricky, since
the values can change in ways that most programmers aren't expecting, although
that integer type is perfectly compatible with Protobuf's wire format. Practical
usage of Pollux as a rule should probably remove \sints{} and \sintl{} from the
\textsc{Int-Int-T} rule. The inclusion of those integer types here also diverges
from Protobuf's official compatibility guide \cite{LanguageGuideProto}.

\subsubsection{Message \& Enum Rules}

\begin{mathpar}
  \infer[Msg-T]{m_1 \preceq m_2}{\mathtt{MSG}\ m_1 \propto \mathtt{MSG}\ m_2}

  \infer[Enum-T]{e_1 \subset e_2}{\enum{e_1} \propto \enum{e_2}}
\end{mathpar}

\subsubsection{Decorator Rules}

Since the type relation is designed for checking for safe type changes, it also
need to ensure that decorator restrictions like moving from a repeated field to
a singleton one aren't violated.

\begin{mathpar}
  \infer[Opt-Add-T]{ }{\tau \propto \optt{\tau}}

  \infer[Opt-Rm-T]{ }{\optt{\tau} \propto \tau}

  \infer[Opt-T]{\tau_1 \propto \tau_2}{\optt{\tau_1} \propto \optt{\tau_2}}

  \infer[Rep-Add-T]{ }{\tau \propto \rept{\tau}}

  \infer[Rep-T]{\tau_1 \propto \tau_2}{\rept{\tau_1} \propto \rept{\tau_2}}
\end{mathpar}

\subsection{Message Relation}~\label{sec:msg-rel}

The descriptor compatibility relation $\preceq$ is defined as below, but first it is
important to discuss the structure of fields and messages.

A message value is a tuple containing a set of reserved field numbers and names,
then a map from an identifier to a tuple with the name of the field and the type
of the field. The identifier is either a tag for a single field, or a set of
integers which is used for a \texttt{oneof}, where multiple fields must be
mutually exclusively set.

Enums will be modeled as a set of integers, likewise ignoring the name given to
each element to favor the integer which is used in the encoding. The enum
relation is defined using $\preceq_e$ to differentiate it from the message
relation.

\subsubsection{Basic Rules}

A compatibility relation is both reflexive and transitive, although it is
neither symmetric nor anti-symmetric.  

\begin{mathpar}
   \infer[Refl-M]{ }{m \preceq m}

   \infer[Refl-E]{ }{e \preceq_e e}

   \infer[Trans-M]{m_1 \preceq m_2 \\ m_2 \preceq m_3}{m_1 \preceq m_3}

   \infer[Trans-E]{e_1 \preceq_e e_2 \\ e_2 \preceq_e e_3}{e_1 \preceq_e e_3}
\end{mathpar}

\subsubsection{Field Update Rules}

For the moment, names are recorded the protobuf field descriptors for accuracy
and to allow the relations to expand in the future. However, since names are not
recorded in the encoded message, they are currently allowed to freely change.

{\color{red} I seem to remember Tej saying that we \emph{should} have the name
  information available somewhere in the descriptor. I need to check my notes
  about that.}

Generally speaking, we can change the type of a field, add a new field or update
a field in a map.

\begin{mathpar}
  \infer[Field-Type]{f(id) = \tau_1 \\ \tau_1 \propto \tau_2}{\msg{r}{f} \preceq \msg{r}{f[id \mapsto \tau_2]}}

  \infer[Field-Add]{id \not \in \textrm{dom}(f) }{\msg{r}{f} \preceq \msg{r}{f[id \mapsto \tau]}}

  % TODO: May need restrictions about maps in maps...
  \infer[Map-Type-F]{f(id) = (\map{\tau_1}{\impt{\tau_2})} \\ \tau_1 \propto \tau_1' \\ \tau_2 \propto \impt{\tau_2'} \\ \tau_1 \ne \byt}
  {\msg{r}{f} \preceq \msg{r}{f[id \mapsto (\map{\tau_1'}{\imp{\tau_2'}})]}}
\end{mathpar}

Messages are similar to records in Lean or other functional languages and
message compatibility can take the same two routes that record subtyping can:
width compatibility and depth compatibility. In width compatibility updates, new
fields are added to the message, making the updated version a strict superset of
the original message. In depth compatibility updates, the number of fields
remains the same, but other attributes of the field are updated compatible ways.

\subsubsection{Oneof Rules}

The \texttt{oneof} rules are a bit complicated because a \texttt{oneof}
struggles against the notion that the field relations only connects one field to
one field by encapsulating several fields into a singular field.

% FIXME: Oneof fields don't cleanly fit into the map abstractions since there is
% no clear choice of key.

\begin{mathpar}
  \infer[Oneof-Intro-Field]{d \in \{\imp, \opt\} \\ f(id) = d\ \tau}
  {\msg{r}{f} \preceq \msg{r}{f[\{id\} \mapsto \oneof{id \mapsto (d\ \tau)}]}}

  % May need some statement about f_{n+1} not already being in the message?
  % You can introduce a new field, but not move an existing one.
  \infer[Oneof-Add-F]{d \in \{\mathtt{IMP}, \mathtt{OPT}\} \\ id_n \not \in
    id_s \\ f(id_s) = m \\ id_n \not \in  \mathrm{dom}(m) \\ \tau \not= \mathtt{MAP}}
  {\msg{r}{f} \preceq \msg{r}{f[id_s \cup \{id_n\} \mapsto (m[id_n \mapsto (d \tau)])]}}

  % This rule should just be a combination of the above two rules...
  \infer[Oneof-Intro-New]{\dom{ f } \cap \dom{ f_o } = \emptyset
    \\ \mathtt{REP} \not \in dec(f_o) \\ \forall\;id_n \in \dom{f_o}.\
    \texttt{snd}\ f_o(id_n) \ne \mathtt{MAP}}{\msg{r}{f} \preceq \msg{r}{f[\dom{f_o}
      \mapsto f_o]}}

  \infer[Oneof-Elim]{f(id_n) = m \\ id \in id_n }
  {\msg{r}{f} \preceq \msg{r}{f[id_n \setminus id \mapsto m \setminus id]}}

  % I'm not sure about the premise of this rule, but it seems better than having
  % to making another version of all the other rules. This is the first instance
  % where having a field relation would be very helpful.
  \infer[Oneof-Field-Update]{f_o(id') = \tau \\ \tau \propto \tau'}
  {\msg{r}{f[id \mapsto f_o]} \preceq \msg{r}{f[id \mapsto ( f_o[id' \mapsto \tau'] )]}}
\end{mathpar}

\subsubsection{Enum Rules}

The enum rules are similar to the message rules at least in the sense that enums
can be updated to have new fields without issue. However, when considering
removing fields from an enum definition, it is important to know if the enum is
treated as open or closed~\cite{EnumBehavior}. By default, \texttt{proto3}
treats all enums as open, so it should be possible to write a rule allowing for
the deletion of an enum field. This naturally aligns with how \texttt{go}
handles enums, since it doesn't support closed enums. Since protobuf enums don't
contain extra information, no depth rules is allowed.

\begin{mathpar}
  \infer[Enum-Width]{vals(e_1) \subset vals(e_2)}{e_1 \preceq_e e_2}
\end{mathpar}

\subsubsection{Reserved Field Rules}

\begin{mathpar}
  \infer[Reserved-Add]{r \subset r'}{\msg{r}{f} \preceq \msg{r'}{f}}

  {\color{red}
    \infer[Reserved-Rm]{r \supset r'}{\msg{r}{f} \preceq \msg{r'}{f}}
  }
\end{mathpar}

\subsubsection{Valid Descriptors}

\begin{definition}[Valid Protobuf Descriptors]
  A protobuf message descriptor is \emph{valid}, denoted $v(d)$ if
  $\msg{\emptyset}{\emptyset} \preceq \msg{r}{f}$.
\end{definition}

\subsubsection{Compatibility Theorem}

The main field-level compatibility theorem is stated below.
\\
\begin{theorem}[Protobuf Compatibility]
  If $d_1 \preceq d_2$ and $d_2$ is well-formed, then for all binary strings
  $bs$ such that $\mathtt{parse}_{d_1}\ bs = \mathtt{Some}\ v_1$,
  $\mathtt{parse}_{d_2}\ bs = \mathtt{Some}\ v_2$ and $v_1 \prec v_2$
\end{theorem}

% Local Variables:
% citar-bibliography: ("../pollux.bib")
% TeX-master: "../pollux.tex"
% jinx-local-words: "protobuf"
% End:
